\documentclass[11pt,a4paper]{article}

% --- Pacchetti Fondamentali ---
\usepackage[utf8]{inputenc}
\usepackage[T1]{fontenc}
\usepackage{lmodern}
\usepackage[italian]{babel}
\usepackage[margin=.75in]{geometry}
\usepackage[table]{xcolor}
\usepackage{enumitem}
\usepackage{fancyhdr}
\usepackage{tabularx}
\usepackage{booktabs}

% --- Configurazione Colori ---
\definecolor{primary}{RGB}{38, 66, 139}

% --- Configurazione Header e Footer ---
\setlength{\headheight}{14pt}
\addtolength{\topmargin}{-2pt}
\pagestyle{fancy}
\fancyhf{}
\fancyhead[L]{\textbf{Nome:} \underline{\hspace{5cm}}}
\fancyhead[C]{\textbf{Data:} \underline{\hspace{2.5cm}}}
\fancyhead[R]{\textbf{Punteggio:} \underline{\hspace{1cm}} / 20}
\fancyfoot[C]{\thepage}
\renewcommand{\headrulewidth}{0pt}

\begin{document}

\begin{center}
    {\Large \textbf{\color{primary} Test Finale di Valutazione}} \\
    \textit{Programmazione e Tecnologie Web}
\end{center}

\vspace{0.5cm}

\begin{enumerate}[label=\textbf{\arabic*.}, leftmargin=1.5em, itemsep=1.5em]

    % Lezione 1: Architettura e HTTP
    \item \textbf{Quale parte di un indirizzo internet (URL) non viene inviata al server?}
    \begin{enumerate}[label=\alph*)]
        \item La Query String (dopo il punto interrogativo).
        \item Il Percorso (Path).
        \item Il Frammento (Fragment, dopo il cancelletto \#).
        \item Il Protocollo (http://).
    \end{enumerate}

    \item \textbf{Quale versione di HTTP ha introdotto le "connessioni persistenti" per essere più veloce?}
    \begin{enumerate}[label=\alph*)]
        \item HTTP/0.9
        \item HTTP/1.0
        \item HTTP/1.1
        \item HTTP/3
    \end{enumerate}

    \item \textbf{Se un'operazione è "Sincrona" e richiede molto tempo, cosa succede alla pagina?}
    \begin{enumerate}[label=\alph*)]
        \item Si chiude da sola.
        \item Si blocca tutto e l'utente non può fare nulla.
        \item Diventa più sicura.
        \item Cambia colore.
    \end{enumerate}

    % Lezione 2: Statico vs Dinamico
    \item \textbf{Cosa fa un generatore di siti statici (come Hugo)?}
    \begin{enumerate}[label=\alph*)]
        \item Crea le pagine ogni volta che un utente le clicca.
        \item Crea tutte le pagine una volta sola, prima di metterle online.
        \item Ha bisogno di un database sempre acceso.
        \item Non permette di usare immagini.
    \end{enumerate}

    \item \textbf{Nel Client-Side Rendering (CSR), chi costruisce la pagina?}
    \begin{enumerate}[label=\alph*)]
        \item Il server invia la pagina già pronta.
        \item Il browser scarica un pezzetto di codice e costruisce tutto da solo.
        \item Il database.
        \item L'utente a mano.
    \end{enumerate}

    \item \textbf{Cosa succede ai dati salvati nel \textit{localStorage} del browser?}
    \begin{enumerate}[label=\alph*)]
        \item Restano lì anche se chiudiamo e riapriamo il browser.
        \item Si cancellano da soli dopo un secondo.
        \item Vengono spediti per posta a casa del programmatore.
        \item Rendono il computer molto più pesante.
    \end{enumerate}

    % Lezione 3: API e JSON
    \item \textbf{In un'architettura a "Microservizi", il sito è fatto da...}
    \begin{enumerate}[label=\alph*)]
        \item Un solo pezzo enorme e difficilissimo da cambiare.
        \item Tanti piccoli programmi separati che si aiutano a vicenda.
        \item Solo disegni fatti a mano.
        \item Robot che scrivono il codice da soli.
    \end{enumerate}

    \item \textbf{Cosa significa "Serializzare" un oggetto in JSON?}
    \begin{enumerate}[label=\alph*)]
        \item Trasformare un dato in una scritta (stringa) per poterlo inviare.
        \item Mettere una password al file.
        \item Inviare tante mail.
        \item Rimpicciolire una foto.
    \end{enumerate}

    \item \textbf{Cosa controlla la fase di "Autorizzazione"?}
    \begin{enumerate}[label=\alph*)]
        \item Chi è l'utente (il suo nome).
        \item Cosa può fare l'utente (i suoi permessi).
        \item Da quale città si collega.
        \item Quanti anni ha.
    \end{enumerate}

    % Lezione 4: Backend
    \item \textbf{Perché diciamo che il Backend (Server) è un posto "sicuro" per i segreti?}
    \begin{enumerate}[label=\alph*)]
        \item Perché l'utente non può vedere cosa succede dentro al server.
        \item Perché i server sono chiusi a chiave in una cassaforte.
        \item Perché il server non usa mai internet.
        \item Perché i programmatori sono tutti amici tra loro.
    \end{enumerate}

    \item \textbf{A cosa serve il comando \texttt{express.json()} in un server Node.js?}
    \begin{enumerate}[label=\alph*)]
        \item A spegnere il server.
        \item A capire i dati in formato JSON che arrivano dal client.
        \item A cambiare il font del sito.
        \item A connettersi al Wi-Fi.
    \end{enumerate}

    \item \textbf{Un metodo HTTP è "Sicuro" (come la \texttt{GET}) quando...}
    \begin{enumerate}[label=\alph*)]
        \item Chiede i dati senza cambiare o rompere nulla sul server.
        \item Blocca i virus e gli hacker.
        \item Fa risparmiare corrente elettrica.
        \item Fa andare il sito molto più veloce.
    \end{enumerate}

    \item \textbf{Perché usiamo \texttt{async} e \texttt{await} in JavaScript?}
    \begin{enumerate}[label=\alph*)]
        \item Per scrivere codice asincrono in modo più semplice e leggibile.
        \item Per rendere il sito più colorato.
        \item Per non usare più le variabili.
        \item Per far funzionare il mouse.
    \end{enumerate}

    % Lezione 5: React
    \item \textbf{In React, a cosa serve \texttt{useState}?}
    \begin{enumerate}[label=\alph*)]
        \item A cambiare il colore del browser.
        \item A memorizzare e aggiornare un dato che cambia nella pagina (es. un contatore).
        \item A inviare una mail.
        \item A spegnere il computer.
    \end{enumerate}

    \item \textbf{Cosa fa un cookie \textit{HttpOnly}?}
    \begin{enumerate}[label=\alph*)]
        \item Si può usare solo se c'è internet.
        \item Impedisce a JavaScript di leggerlo, proteggendo l'utente.
        \item Scade dopo un secondo.
        \item Cancella la cronologia.
    \end{enumerate}

    % Lezione 6: Real-time & Tooling
    \item \textbf{Cosa permettono di fare i WebSocket?}
    \begin{enumerate}[label=\alph*)]
        \item Scaricare file più velocemente.
        \item Creare una connessione sempre aperta per scambiare dati subito (come una chat).
        \item Spegnere il server.
        \item Stampare la pagina.
    \end{enumerate}

    \item \textbf{A cosa serve il comando \texttt{npm install}?}
    \begin{enumerate}[label=\alph*)]
        \item A installare le librerie necessarie al progetto.
        \item A installare Windows.
        \item A cancellare tutto il codice.
        \item A cambiare il nome del programmatore.
    \end{enumerate}

    \item \textbf{In React, perché si usa \texttt{useCallback}?}
    \begin{enumerate}[label=\alph*)]
        \item Per scaricare dati dal server.
        \item Per non ricreare una funzione inutilmente ogni volta, risparmiando fatica.
        \item Per fare il login.
        \item Per ingrandire le immagini.
    \end{enumerate}

    % REST & Async
    \item \textbf{Cosa restituisce il comando \texttt{fetch()}?}
    \begin{enumerate}[label=\alph*)]
        \item Subito il risultato finale.
        \item Una Promise (una promessa di un dato che arriverà).
        \item Un errore.
        \item Una scritta HTML.
    \end{enumerate}

    \item \textbf{Quale metodo HTTP si usa per sostituire completamente un dato esistente?}
    \begin{enumerate}[label=\alph*)]
        \item \texttt{POST}
        \item \texttt{GET}
        \item \texttt{PUT}
        \item \texttt{DELETE}
    \end{enumerate}

\end{enumerate}

\vspace{1cm}
\begin{center}
\begin{tabular}{|l|l|}
\hline
\textbf{Punteggio Totale (su 20)} & \hspace{3cm} / 20 \\ \hline
\textbf{Valutazione Finale} & \hspace{3cm} \\ \hline
\end{tabular}
\end{center}

\newpage
\fancyhead[L]{Griglia di Valutazione}
\fancyhead[R]{Riservato al Docente}

\vspace*{1cm}

\begin{center}
    {\Large \textbf{Griglia di Valutazione e Correzione}}
\end{center}

\vspace{1cm}

\begin{table}[ht]
\centering
\begin{tabularx}{\textwidth}{@{}|l|X|l|@{}}
\hline
\rowcolor{gray!20} \textbf{Domanda} & \textbf{Argomento} & \textbf{Risposta Corretta} \\ \hline
1 & Struttura URL (Frammento) & c \\ \hline
2 & Evoluzione HTTP (1.1) & c \\ \hline
3 & Sincronia e Blocco UI & b \\ \hline
4 & SSG (Hugo) & b \\ \hline
5 & Client-Side Rendering & b \\ \hline
6 & LocalStorage (Persistenza) & a \\ \hline
7 & Microservizi (Struttura) & b \\ \hline
8 & Serializzazione JSON & a \\ \hline
9 & Autorizzazione (Permessi) & b \\ \hline
10 & Ambiente Sicuro (Server) & a \\ \hline
11 & Middleware JSON & b \\ \hline
12 & Metodi Sicuri (GET) & a \\ \hline
13 & Async/Await & a \\ \hline
14 & React useState & b \\ \hline
15 & Cookie HttpOnly & b \\ \hline
16 & WebSocket & b \\ \hline
17 & Tooling (npm install) & a \\ \hline
18 & React useCallback & b \\ \hline
19 & Fetch API (Promise) & b \\ \hline
20 & Metodo PUT & c \\ \hline
\end{tabularx}
\end{table}

\vspace{1cm}
\textbf{Criteri di Valutazione:}
\begin{itemize}
    \item 1 punto per ogni risposta corretta.
    \item 0 punti per risposta errata o non data.
    \item \textbf{Soglia di sufficienza:} 12/20 (60\%).
\end{itemize}

\end{document}
